\section{Future Modules}

The architecture is intended to support additional benchmark modules without duplicating the central planning and materialization logic.
Shared functionality such as configuration loading, reference resolution, force-field conversion, execution, and evaluation should remain in \texttt{pipeline/}, \texttt{converter/}, \texttt{resolvers/}, and \texttt{parsers/}.
Module-specific code should only describe module-specific scientific workflows.

\subsection{Module B: Mixture Adsorption and Selectivity}

Module B will extend the current single-component workflow to mixtures such as CO2/N2.
Expected additions include:

\begin{itemize}
    \item mixture composition in the benchmark configuration,
    \item multiple molecule templates,
    \item mixture GCMC input generation,
    \item component-wise adsorption summaries,
    \item selectivity calculation,
    \item reference comparison for mixture data where available.
\end{itemize}

\subsection{Module C: Widom Insertion and Henry Coefficients}

Module C will focus on low-pressure adsorption properties and potential comparison.
Expected additions include:

\begin{itemize}
    \item Widom insertion input generation,
    \item Henry coefficient extraction,
    \item optional force-field comparison,
    \item low-pressure reference matching,
    \item uncertainty estimation for insertion averages.
\end{itemize}

\subsection{Module D: Framework Flexibility}

Module D will address flexible-framework simulations.
Expected additions include:

\begin{itemize}
    \item flexible framework preparation,
    \item MD or hybrid MD/GCMC workflows,
    \item trajectory analysis,
    \item OVITO-based visualization outputs,
    \item comparison between rigid and flexible adsorption behavior.
\end{itemize}

\subsection{Avoiding Redundancy}

Preparation and materialization should remain central whenever possible.
Individual modules should not duplicate generic logic such as copying files, creating run directories, resolving references, or writing evaluation outputs.
Instead, modules should provide the module-specific parts of the run plan that the central pipeline can materialize and execute.
